\chapter{Observable Lists and Large Data}
\label{ch:observable-lists}
\chapterquote{A large table is usually not slow because it is large; it is slow because the program lies about what changed.}{A table performance rule}

\begin{chaptermap}
Core question & How should changing row collections be represented before Swing displays them?\\
Primary APIs & \api{ObservableList}, \api{ObservableArrayList}, \api{ListChange}, \api{ListChangeSet}, \api{FilteredList}, \api{SortedList}, \api{MappedList}, \api{LoadableList}, \api{CollectionItemValue}, \api{CollectionMatchingValue}, \api{CollectionDerivedValue}, \api{MasterDetailCollection}, \api{EdtObservableList}, and \api{GlazedObservableList}.\\
Plain Swing baseline & \api{AbstractTableModel}, \api{ListModel}, and \api{ComboBoxModel} need precise events but are often fed broad resets.\\
Corusco connection & Observable lists adapt to list, combo-box, and table models with precise events and lifecycle-owned bindings.\\
Companion example & \largeDataExample.\\
\end{chaptermap}

\section{A list used by a screen is not just storage}

A Java \api{List} answers questions about current contents: how many elements
exist, which element is at index five, and what the current snapshot looks like.
A Swing screen needs more. It needs to know what changed. Did a row appear? Did
a row disappear? Was one row replaced? Did the order change? Was the entire
collection cleared? Should the table repaint one cell, five rows, or the whole
viewport?

\termdef{Observable list}{in Corusco} is an indexed collection that publishes
ordered structural changes. The core contract includes reads, mutations,
batches, immutable snapshots, and subscriptions. It is the collection sibling of
\api{ReadableValue}.

\begin{lstlisting}[caption={The shape of the observable-list contract.}]
ObservableList<CustomerRow> rows = ObservableArrayList.empty();

Subscription subscription = rows.subscribe(changes -> {
    for (ListChange<CustomerRow> change : changes.changes()) {
        System.out.println(change);
    }
});

rows.add(new CustomerRow("Ada", "active"));
rows.set(0, new CustomerRow("Ada", "preferred"));
subscription.close();
\end{lstlisting}

The list is still an ordinary indexed data source. The important addition is
that mutations are observable as facts, not merely as "something changed".

\begin{principlebox}{Tell the Truth About Change}
Large-data Swing programming begins with honest events. If one row changed,
report one row. If a batch inserted ten rows, report the inserted range. Use a
reset only when the program really cannot describe the difference.
\end{principlebox}

\begin{figure}[htbp]
\centering
\begin{minipage}[t]{0.48\linewidth}
\includegraphics[width=\linewidth]{\largeDataFigure}
\end{minipage}\hfill
\begin{minipage}[t]{0.48\linewidth}
\includegraphics[width=\linewidth]{\tableCoordinatesFigure}
\end{minipage}
\caption{Observable row changes become visible Swing behavior. A large table needs precise data events and a clear distinction between model coordinates, view coordinates, and row identity.}
\label{fig:observable-list-screens}
\end{figure}

\textbf{Why broad resets hurt.}

Many table models solve every update by calling \api{fireTableDataChanged}. That
method is a useful escape hatch, but a poor default. It tells the table that the
model's contents may have changed broadly. The table must invalidate more state
than necessary. Selection may need repair. Sorters and row filters may do extra
work. Renderers repaint cells whose values did not change. Users see flicker or
lost context.

For ten rows, the bad habit is invisible. For ten thousand rows, it becomes a
user-interface problem. For a table with sorting, filtering, column persistence,
cell validation, and selection preservation, it becomes a correctness problem.

\begin{lstlisting}[caption={The common but imprecise table-model update.}]
void replaceCustomer(CustomerRow updated) {
    int index = findById(updated.id());
    rows.set(index, updated);
    fireTableDataChanged();       // The model knows more than it reports.
}
\end{lstlisting}

The model knows that one row was replaced. It should say so.

\begin{lstlisting}[caption={A precise table-model update.}]
void replaceCustomer(CustomerRow updated) {
    int index = findById(updated.id());
    rows.set(index, updated);
    fireTableRowsUpdated(index, index);
}
\end{lstlisting}

Corusco's observable lists preserve this idea before Swing enters the story.
The presenter can replace row five; the Swing adapter can translate that into
the proper table event.

\textbf{The core list contract.}

\api{ObservableList<E>} exposes size, indexed read, immutable snapshot, add,
insert, set, remove, move, clear, batch, and subscribe. The operations are
deliberately explicit. Mutating through a hidden iterator would make change
reporting ambiguous. The list owns the mutation vocabulary.

\begin{lstlisting}[caption={Core mutation vocabulary.}]
rows.add(element);
rows.add(index, element);
rows.set(index, replacement);
rows.remove(index);
rows.move(fromIndex, toIndex);
rows.clear();
rows.batch(list -> {
    list.add(first);
    list.add(second);
});
\end{lstlisting}

A batch groups inner mutations into one delivered change set. The order of
changes still matters. A subscriber can see "insert row, insert row, replace
row" as one delivery and update a dependent structure once.

\textbf{List changes and change sets.}

\termdef{List change}{in Corusco} is one structural mutation: inserted,
removed, replaced, moved, or cleared. \termdef{List change set}{in Corusco} is
an ordered delivery of one or more changes. Batches use change sets to avoid
publishing partial intermediate states. Each delivered change set also carries a
\api{ChangeOrigin}, using \api{StandardChangeOrigin.MODEL} for ordinary
presenter-owned writes unless the caller supplies a more specific origin.

This is more precise than a plain "changed" callback and less elaborate than a
complete diff engine. It matches what Swing adapters normally need. A table
model can turn an inserted change into \api{fireTableRowsInserted}; a removed
change into \api{fireTableRowsDeleted}; a replaced change into
\api{fireTableRowsUpdated}; a clear into a deletion range or reset according to
the model's state.

\begin{pitfallbox}{Indexes Are Temporal}
Indexes inside a change describe the list at a particular point in the ordered
change stream. Do not collect indexes from several changes and apply them later
without respecting order.
\end{pitfallbox}

\textbf{ObservableArrayList.}

\api{ObservableArrayList} is the ordinary storage implementation. Use it when
the presenter owns the rows directly. It is not a database, not a virtualized
data source, and not a replacement for paging. It is a precise in-memory row
source.

\begin{lstlisting}[caption={Presenter-owned row storage.}]
final class CustomerRows {
    private final ObservableArrayList<CustomerRow> rows =
            ObservableArrayList.empty();

    ObservableList<CustomerRow> rows() {
        return rows;
    }

    void loaded(List<CustomerRow> loadedRows) {
        rows.batch(list -> {
            list.clear();
            loadedRows.forEach(list::add);
        });
    }
}
\end{lstlisting}

This example still clears and repopulates on load. That may be acceptable for a
fresh search result. It is not ideal for a refresh where row identity should be
preserved. The point is that even broad replacement is named by the row source,
not hidden in the table.

\textbf{Filtering.}

\api{FilteredList<E>} exposes a read-only view over a source list using a
predicate. The source remains the mutation owner. The filtered view subscribes
to source changes and reports changes in the filtered coordinate space.

\begin{lstlisting}[caption={Filtered view over a source list.}]
ObservableArrayList<CustomerRow> source = ObservableArrayList.empty();
FilteredList<CustomerRow> activeRows =
        FilteredList.of(source, row -> row.status() == CustomerStatus.ACTIVE);
\end{lstlisting}

Filtering is easy to write badly in Swing. A renderer decides whether to gray
out rows. A row filter reads values through view indexes. A search listener
rebuilds the entire table model on each keystroke. Corusco's core view keeps
the filter before Swing and makes the relationship testable.

\begin{migrationbox}{Filter Before Rendering}
If a renderer contains business filtering logic, move that logic into a filtered
row source or table row sorter policy. Renderers should present visible rows,
not decide which rows exist.
\end{migrationbox}

\section{Sorting}

\api{SortedList<E>} exposes a sorted read-only view. Swing also has
\api{RowSorter}, and it is often useful. The choice depends on ownership. If the
sort order is purely a table-view preference, a Swing row sorter is appropriate.
If a presenter-owned downstream model needs sorted order before Swing, a sorted
observable list is appropriate.

The distinction matters for selection. Swing view indexes change when the table
sorter changes. Model indexes refer to the table model. A presenter-owned sorted
list is itself the model order. Do not mix these coordinate systems casually.

\begin{lstlisting}[caption={Sorted view for presenter-owned order.}]
SortedList<CustomerRow> byName = SortedList.of(source,
        Comparator.comparing(CustomerRow::displayName)
                .thenComparing(CustomerRow::id));
\end{lstlisting}

\textbf{Mapping.}

\api{MappedList<S,T>} converts each source element into a view element. This is
useful for combo-box labels, small summary lists, and table helper projections.
The source remains the mutation owner. The mapped list is a view.

\begin{lstlisting}[caption={Mapping rows to display choices.}]
MappedList<CustomerRow, CustomerChoice> choices = MappedList.of(source,
        row -> new CustomerChoice(row.id(), row.displayName()));
\end{lstlisting}

Do not use mapping to hide expensive per-row work that should be cached or
loaded separately. Mapping should be cheap enough to run as part of change
propagation.

\textbf{Streams and snapshots.}

\api{ObservableReadableCollection.stream()} is for immediate read-only
traversal in collection order. It is useful for questions such as "does any row
match this id?" or "how many visible rows satisfy this predicate?" without
forcing an immutable snapshot copy. The returned stream is still an ordinary
single-use Java stream, and it is not a durable snapshot.

\api{snapshot()} and \api{snapshotStream()} are the durable alternatives. Use
them when the contents must survive later list mutation, be stored in a value,
or be passed across a boundary that may outlive the current read.

\begin{lstlisting}[caption={Immediate traversal versus durable snapshot.}]
boolean hasErrors = rows.stream()
        .anyMatch(CustomerRow::hasErrors);

List<CustomerRow> retainedRows = rows.snapshot();
long retainedErrorCount = rows.snapshotStream()
        .filter(CustomerRow::hasErrors)
        .count();
\end{lstlisting}

\textbf{Loadable lists.}

\api{LoadableList<E>} models a collection that is loaded through a supplier and
can be detached. It is useful for presenter-owned cached row sets where the
screen should release cached rows without discarding the row-source object.

\begin{lstlisting}[caption={A detachable loaded row source.}]
LoadableList<CustomerRow> rows = LoadableList.of(repository::loadCustomers);

List<CustomerRow> firstSnapshot = rows.snapshot();
rows.detach();
List<CustomerRow> secondSnapshot = rows.snapshot();
\end{lstlisting}

As with \api{LoadableValue}, this does not make slow loading safe for the EDT.
If loading may block, put the load behind a task boundary and update the list on
the appropriate thread. The loadable list models storage and detachment, not UI
threading.

\textbf{EDT adaptation.}

Core observable lists are not Swing dispatchers. Listener delivery is
synchronous on the mutating thread. Swing table models, list models, and combo
box models are EDT-confined. The boundary must be explicit.

\api{EdtObservableList} wraps a source list and delivers wrapper subscriber
callbacks on the EDT. If the source fires on the EDT, delivery is immediate. If
the source fires elsewhere, delivery is queued with \api{SwingUtilities.invokeLater}.
The wrapper does not make the source list itself thread-safe.

\begin{lstlisting}[caption={Adapting a row source before Swing consumes it.}]
ObservableList<CustomerRow> source = presenter.rows();
EdtObservableList<CustomerRow> swingRows = EdtObservableList.of(source);

ObservableTableModel<CustomerRow> model =
        ObservableTableModel.of(swingRows,
                GeneratedCustomerRowTableDescriptor.DESCRIPTOR);
\end{lstlisting}

\begin{pitfallbox}{EDT adapter is not a lock}
\api{EdtObservableList} changes callback delivery for subscribers of the
wrapper. It does not make arbitrary concurrent mutations of the wrapped source
safe. Use a coherent source ownership policy.
\end{pitfallbox}

\textbf{ObservableTableModel.}

\api{ObservableTableModel<R>} is the direct Swing table model adapter. It reads
cell values through a typed \api{TableDescriptor}. It subscribes to an
\api{ObservableList}. It translates list changes into table events. It is
EDT-confined and closes its source-list subscription when closed.

This combination addresses several common Swing table failures:

\begin{itemize}
\item Row storage is not owned by \api{JTable}.
\item Column meaning is not inferred from localized headers.
\item Cell values are read through typed column descriptors.
\item Immutable record rows can be edited by replacing the row.
\item Source-list changes become precise table events where possible.
\item Closing the table model removes the source subscription.
\end{itemize}

\begin{lstlisting}[caption={Descriptor-backed observable table model.}]
ObservableTableModel<CustomerRow> model =
        ObservableTableModel.of(rows,
                GeneratedCustomerRowTableDescriptor.DESCRIPTOR);

table.setModel(model);
\end{lstlisting}

\textbf{Editing immutable rows.}

Traditional table models often mutate bean properties directly. That is natural
for mutable beans but awkward for records. Corusco table columns can update a
row by returning a replacement row. The observable list then reports a row
replacement, and the table model fires the appropriate cell or row update.

\begin{lstlisting}[caption={Conceptual immutable-row edit.}]
Column<CustomerRow, String> nameColumn = Column.of(
        CustomerColumns.NAME,
        CustomerRow::displayName,
        (row, value) -> row.withDisplayName(value));
\end{lstlisting}

The exact generated code may differ, but the ownership is the same. The table
does not own the row's fields. The descriptor knows how to read and update the
row. The list owns row replacement.

\section{Selection and row identity}

Large data sets need stable row identity. A selected view row is not enough.
After sorting, filtering, insertion, or refresh, the same view index may refer
to a different row. Selection should be preserved by row identity when the user
expects continuity.

Corusco table selection binding connects Swing selection to typed model state:
selected model row, selected row object, or selected row id according to the
screen's needs. The observable list chapter matters because precise row changes
give the binding a fighting chance to preserve context.

\textbf{GlazedLists interop.}

GlazedLists is a mature library for event lists, filtering, sorting, and
Swing-table integration. Corusco does not need to pretend that this experience
does not exist. \api{GlazedObservableList} adapts a GlazedLists \api{EventList}
to the Corusco \api{ObservableList} contract.

Use this when an application already depends on GlazedLists or needs its
advanced list machinery. The boundary remains important: decide which library
owns filtering, sorting, threading, and disposal. Adapt deliberately.

\begin{coruscobox}{Interop Instead of Replacement}
Corusco's observable-list contract is a presenter and binding contract. It can
own storage directly or adapt a mature external event-list source such as
GlazedLists. The goal is not to force every application through one collection
engine.
\end{coruscobox}

\textbf{Large-data discipline.}

Large data work starts with a simple question: can the user see all of this at
once? If not, avoid doing component-level work for every row. A table with
100000 rows should not allocate 100000 components. Swing renderers already help
by stamping visible cells. Your model must help by avoiding unnecessary resets,
expensive per-cell computation, and broad repaint storms.

Useful rules:

\begin{itemize}
\item Keep row objects compact and meaningful.
\item Make renderers cheap and complete.
\item Keep filtering and sorting out of renderers.
\item Preserve row identity across refresh.
\item Batch bulk changes.
\item Avoid firing one event per row when one ordered change set can describe
the batch.
\item Measure with realistic row counts, not a five-row demonstration.
\end{itemize}

\textbf{Plain Swing baseline.}

A careful plain Swing model can be perfectly acceptable:

\begin{lstlisting}[caption={Handwritten precise table model.}]
final class CustomerTableModel extends AbstractTableModel {
    private final List<CustomerRow> rows = new ArrayList<>();

    void addCustomer(CustomerRow row) {
        int index = rows.size();
        rows.add(row);
        fireTableRowsInserted(index, index);
    }

    void replaceCustomer(int index, CustomerRow row) {
        rows.set(index, row);
        fireTableRowsUpdated(index, index);
    }

    @Override
    public int getRowCount() {
        return rows.size();
    }
}
\end{lstlisting}

This code is honest. Its limitation is reuse. Every table model must repeat row
storage, event translation, column metadata, editing policy, resource identity,
and cleanup. Corusco extracts those repeated concerns.

\textbf{Corusco treatment.}

The Corusco version splits ownership:

\begin{itemize}
\item The presenter owns an \api{ObservableList<CustomerRow>}.
\item The descriptor owns columns and cell access.
\item The Swing table model owns translation to \api{TableModelEvent}.
\item The binding or scope owns cleanup.
\item The table owns viewport, selection UI, renderers, and editors.
\end{itemize}

This is not more abstract for its own sake. It gives each recurring fact a
stable owner. The same row source can be tested without Swing, displayed in a
table, adapted to a list, or used by a generated table companion.

\textbf{Worked example.}

\begin{examplebox}{Large Data}
\largeDataExample{} builds a deterministic row set, shows it through a table,
and separates row generation from Swing display. Use it as a base for
screenshot-scale tables and for model-level tests that run without a visible
window.
\end{examplebox}

Read the example by following row identity. Where are rows created? Which code
owns the list? Which code owns the table model? Which code can be tested without
Swing? If these answers are visible, the screen can grow.

\textbf{Review notes.}

Observable lists do not remove the need for paging, server-side filtering, or
database indexes. They make in-memory presentation changes honest. A million-row
database result should not be blindly loaded into an observable array list just
because the UI can subscribe to it. Use the right data strategy first; then
publish the presentation subset precisely.

The most useful design habit is to draw the boundary between the authoritative
collection and the presentation collection. In a small utility, those two may be
the same object: an observable array list that directly stores the visible rows.
In a business system, the authoritative data may live in a repository, a cache,
a server-side cursor, or a GlazedLists event list. The presentation list is then
only the current Swing-facing projection. Once this boundary is named, questions
about filtering, sorting, selection, paging, and disposal become easier because
they are no longer hidden inside the table.

Row identity should be chosen before row order. A table refresh is not merely a
new sequence of objects; it is an opportunity to preserve the user's place. If a
customer row has a stable customer id, selection can survive sorting, filtering,
and replacement. If a row is identified only by its current integer position,
every refresh destroys meaning. This is why descriptors and row models should
expose business identity where the screen needs continuity, even when the table
itself continues to paint ordinary row indexes.

The precision of change events is not an optimization to postpone until the end.
It shapes the user's trust in the screen. A broad reset tells a table that many
assumptions are gone, so sorters, selections, editors, and repaint decisions may
all become conservative. A precise insert, update, or remove says which facts
changed and lets the rest of the view stay still. The user experiences this as a
table that does not jump unnecessarily.

Batching deserves similar respect. Firing one hundred single-row events can be
correct and still unpleasant if observers perform per-event work. Firing one
opaque reset can be fast and still destructive. A good observable-list API makes
the middle path available: related mutations are grouped, but the group still
describes what happened. Corusco's collection support is valuable exactly
because it keeps that discipline close to ordinary row operations.

Threading must be decided at the list boundary, not at random subscribers. A
repository callback may arrive on a pool thread. A GlazedLists event may arrive
under its own locking policy. A Swing table model must be touched on the EDT.
If every subscriber tries to repair threading independently, the program becomes
unreviewable. Wrap the source, adapt the event stream, or publish into an
EDT-confined presentation list before Swing sees the changes.

Filtering and sorting are also ownership decisions. A view-local sort installed
by \api{JTable} is useful when ordering is a visual preference and the domain
model does not care. A presenter-owned sorted list is better when other screen
facts depend on the same order, when the order must be tested without Swing, or
when several components share the sorted projection. The important point is to
avoid having two different owners silently sort the same rows in different ways.

Mapping is attractive because it can keep rows small and presentation-specific.
It is dangerous when it hides expensive work. A mapped list that formats dates
or computes small value objects is ordinary. A mapped list that performs service
calls, loads icons from disk, or consults remote permissions is a latency trap.
The map function should be cheap, deterministic, and free of Swing component
state; otherwise the list becomes a disguised background task.

Memory use is part of list design. A screen that shows a current work queue may
comfortably keep every visible row in memory. A screen that searches years of
audit history probably should not. Observable lists describe presentation
mutation; they do not require the whole domain universe to be present. Use
server-side paging, incremental loading, or summary/detail loading when the
business problem is larger than the screen.

The companion examples should therefore be read as scale models, not as data
strategy commandments. The large-data example creates deterministic rows so
screenshots and tests are stable. A production screen may replace that source
with a repository-backed loadable list or with a GlazedLists adapter. The table,
descriptor, renderer, and selection lessons remain the same because they depend
on the presentation contract rather than the particular data source.

When reviewing a table screen, ask for a short story of one row. Where is it
created? How does it enter the observable list? How does a cell value reach the
renderer? What event is fired when the row changes? How is selection preserved
when a new version of the row arrives? Where are subscriptions closed? A design
that can answer those questions is usually ready for more rows.

The same story should be told for removal. A row that disappears may be filtered
out, deleted from the underlying data set, moved to another page, or hidden by a
permission change. The observable event is similar, but the user meaning is not.
If the screen must explain the disappearance, preserve enough model context to
produce a status message or undo action instead of treating every removal as a
silent table mutation.

For very large data, consider whether the observable list should represent all
known rows or only the visible window of rows. A virtualized or paged source may
publish a compact presentation slice while the repository owns the full result
set. Swing still paints a table, but the table no longer implies that every
possible row is an in-memory object. This is often the difference between a
responsive business screen and an impressive demo that cannot survive real data.

\section{Change Sets, Coordinates, and Derived Views}

The list event model is worth studying because it is the bridge between domain
rows and Swing's strict model contracts. A \api{ListChange.Inserted} says that
contiguous elements were inserted at one index. A \api{ListChange.Removed} says
which contiguous elements occupied an index before removal. A
\api{ListChange.Replaced} says one element was replaced at one index. A
\api{ListChange.Moved} says one element moved from one index to another. A
\api{ListChange.Cleared} says the list removed all elements and carries the old
elements. These are small records, but they encode enough truth for table, list,
combo-box, filtered, mapped, and diagnostic adapters to do precise work.

The indices are in the coordinate system of the list that emits the change. This
sentence is more important than it first appears. A source list emits source
coordinates. A filtered view emits filtered-view coordinates. A sorted view
emits sorted-view coordinates. A mapped list emits mapped-view coordinates that
usually correspond one-to-one with source positions. Swing table view rows are a
different coordinate system again when a row sorter is installed. Most large
table bugs are coordinate bugs wearing performance clothes.

\begin{lstlisting}[caption={A change is interpreted in the emitting list's coordinates.}]
ObservableArrayList<CustomerRow> source = ObservableArrayList.empty();
FilteredList<CustomerRow> active = FilteredList.of(source,
        row -> row.status() == CustomerStatus.ACTIVE);

active.subscribe(changes -> {
    for (ListChange<CustomerRow> change : changes.changes()) {
        // The index belongs to the filtered view, not the source list.
        updateActiveRowsDiagnostics(change);
    }
});
\end{lstlisting}

This rule prevents a tempting shortcut. If a filtered list reports that one
visible row was removed at index three, that does not mean source index three
was removed. The filtered view has already translated the source event through
its predicate. A Swing adapter over the filtered view should fire Swing events
in filtered coordinates. A presenter that needs source identity should store row
ids rather than trying to reverse-engineer source indexes.

Observable collections also bridge back to scalar values. A selected table index
can become a nullable current-row value with \api{CollectionItemValue}. A
selected id can become a nullable current-row value with
\api{CollectionMatchingValue}. Row count, empty state, and snapshot summaries
can be represented with \api{CollectionDerivedValue}. Use its stream-based
factories, such as \api{fromStream}, \api{anyMatch}, \api{allMatch}, and
\api{noneMatch}, when the scalar only needs immediate traversal. Use the
snapshot-based factory when the mapper must retain or inspect a durable
point-in-time list. A selected master row can choose the active detail rows
through \api{MasterDetailCollection}. These bridges preserve change origins, so
a user-originated selection and a model-originated reload remain
distinguishable after crossing the collection boundary.

Batches preserve ordered change semantics while reducing delivery noise. A batch
does not mean "some unknown reset happened." It means a set of mutations belongs
together for delivery. The inner order still matters because later indices are
interpreted after earlier mutations. This is why subscribers should process
changes in the delivered order. Sorting the changes by index or saving them for
later can corrupt the meaning.

\begin{lstlisting}[caption={Process a change set in order.}]
void listChanged(ListChangeSet<CustomerRow> changeSet) {
    for (ListChange<CustomerRow> change : changeSet.changes()) {
        switch (change) {
            case ListChange.Inserted<CustomerRow> inserted -> insert(inserted);
            case ListChange.Removed<CustomerRow> removed -> remove(removed);
            case ListChange.Replaced<CustomerRow> replaced -> replace(replaced);
            case ListChange.Moved<CustomerRow> moved -> move(moved);
            case ListChange.Cleared<CustomerRow> cleared -> clear(cleared);
        }
    }
}
\end{lstlisting}

This processing style also makes tests clearer. A test can assert that a batch
contains three changes in a specific order. It can assert that a filtered view
emits no event when a source insertion does not satisfy the predicate. It can
assert that replacing an inactive row with an active row becomes an insertion in
the filtered view. These are model tests, not Swing tests, and they catch the
hard part before table rendering enters the picture.

Snapshots are deliberately immutable. A snapshot lets a caller inspect current
contents without receiving the mutable storage object. This matters for
diagnostics, tests, and background preparation. If a presenter hands out the
internal list and another object mutates it directly, the observable-list
contract is broken. Mutations must pass through the list's explicit methods so
events can remain truthful.

\begin{pitfallbox}{Do Not Mutate Behind the Observable List}
If some code keeps a reference to the backing collection and modifies it
directly, subscribers do not receive a truthful change set. Treat the observable
list's mutation methods as the only public mutation path.
\end{pitfallbox}

Element immutability is a separate question. Replacing a row object is
observable. Mutating a field inside a row object is not, unless the row object
itself has observable fields and the table model knows about them. For ordinary
table rows, immutable records with replacement updates are easier to reason
about. A cell edit creates a replacement row; the list emits a replacement; the
table model fires an update. The row's old and new values are visible in the
event stream.

\begin{lstlisting}[caption={Immutable row edit through replacement.}]
CustomerRow old = rows.get(modelRow);
CustomerRow updated = old.withStatus(CustomerStatus.PREFERRED);
rows.set(modelRow, updated);
\end{lstlisting}

Mutable rows can still be used, but they need a policy. If a row object changes
internally and remains equal to itself, an observable list cannot infer which
cell changed. The table model may need a row-level notification method, or the
row object may need its own observable fields. Do not accidentally mix mutable
domain objects with a list event model that assumes replacement.

Filtered views are read-only because source ownership matters. A filtered view
can expose rows accepted by a predicate and translate source changes into
visible changes. If callers could insert directly into the filtered view, the
API would have to decide where the element belongs in the source list, whether it
must satisfy the predicate, and what happens when it does not. Those policies
belong to the source owner. The filtered view observes and translates.

The filter predicate can change. The implementation reports predicate
replacement as a reset of the filtered view: previous visible contents are
removed and new visible contents are inserted where applicable. That is a
conservative and deterministic contract. It avoids pretending that a predicate
change is a sequence of ordinary source mutations. A user changing the search
text has changed the view rule, not edited every row.

\begin{lstlisting}[caption={Changing the filter policy resets the filtered view.}]
FilteredList<CustomerRow> filtered = FilteredList.of(source,
        row -> row.status() == CustomerStatus.ACTIVE);

filtered.setPredicate(row -> row.displayName().contains(query.value()));
\end{lstlisting}

Sorted views are also read-only, and their event policy is deliberately
conservative. Sorting makes many source mutations ambiguous to translate as one
precise insert, remove, replace, or move in sorted coordinates. A replacement
may change both row contents and sorted position. A comparator change may
reorder everything. Corusco's sorted view therefore reports visible content
changes as reset-style changes: clear old contents, insert current contents
where needed. This is honest. It may be less granular than a hand-optimized sort
adapter, but it is deterministic and safe.

Use a sorted observable view when order is part of presenter state. Use Swing's
\api{RowSorter} when sort order is a local table-view preference. Use GlazedLists
or another mature event-list engine when the application needs advanced sorted
and filtered event machinery with stronger incremental behavior. Corusco's
contract is not a claim that one sorted implementation is best for every data
volume. It is a way to keep ownership visible.

Mapped views are usually the most straightforward. A source row becomes a display
row, choice object, key-label pair, or small presentation record. Insertions map
to insertions, removals map to removals, replacements map to replacements, moves
map to moves, and clears map to clears. The mapper is invoked when the view is
created and when source changes arrive. The mapper should be cheap,
deterministic, and free of Swing component state.

\begin{lstlisting}[caption={Map domain rows into stable combo choices.}]
record CustomerChoice(CustomerId id, String label) {
    @Override
    public String toString() {
        return label;
    }
}

MappedList<CustomerRow, CustomerChoice> choices = MappedList.of(rows,
        row -> new CustomerChoice(row.id(), row.displayName()));
\end{lstlisting}

This pattern is common with combo boxes. The combo box can display a
\api{CustomerChoice} while the presenter keeps stable ids. The visible string is
not the domain state; it is the rendered label for a choice. If localization or
formatting changes, the mapped view can be rebuilt or the source rows can be
replaced. The model remains typed.

Loadable lists are useful when rows have a reusable owner but an attachable
cache. The first read or mutation loads the cache through a supplier. Detach or
invalidate releases cached rows without firing a list change. Refresh reloads
and emits a reset-style change set only when a previously attached snapshot
changes. As with loadable values, the loader runs on the caller's thread. This
is a storage and lifecycle helper, not a background worker.

\begin{lstlisting}[caption={Loadable rows have cache lifecycle, not magic threading.}]
LoadableList<CustomerRow> rows = LoadableList.of(repository::loadCustomers);

List<CustomerRow> first = rows.snapshot(); // invokes loader
rows.detach();                             // releases cache silently
List<CustomerRow> next = rows.refresh();   // reloads and may notify
\end{lstlisting}

Detach without event is correct when the view is not supposed to observe a user
visible change. A hidden tab releasing its row cache does not mean the user
deleted rows. Refresh is the explicit operation for visible reload. If the
screen needs to show "detached", "loading", "failed", or "loaded", model that as
a separate state value or task state. Do not overload list contents to carry
every lifecycle fact.

Derived views have lifetimes. A filtered list subscribes to its source. A sorted
list subscribes to its source. A mapped list subscribes to its source. A
GlazedLists adapter subscribes to its source. A Swing model subscribes to its
source. Each relationship should be closed when the owner no longer needs it.
If a temporary filtered view is created for a dialog and never closed, the source
list may keep the dialog's graph reachable.

\section{Swing Adapters, Threading, and Lifecycle}

Swing's list and table models are strict about thread ownership. They must be
created and used on the EDT. Corusco's core observable lists are Swing-free and
dispatch synchronously on the mutating thread. The boundary is not optional. A
source list can be mutated only on the EDT while Swing is subscribed, or it can
be wrapped with \api{EdtObservableList} so subscribers of the wrapper receive
callbacks on the EDT. The wrapper does not make the source itself thread-safe.

\begin{lstlisting}[caption={A background source still needs a coherent ownership policy.}]
ObservableArrayList<CustomerRow> source = ObservableArrayList.empty();
EdtObservableList<CustomerRow> edtRows = EdtObservableList.of(source);

SwingUtilities.invokeLater(() -> {
    ObservableTableModel<CustomerRow> model =
            ObservableTableModel.of(edtRows, customerTableDescriptor);
    table.setModel(model);
});
\end{lstlisting}

This listing adapts callback delivery, not arbitrary concurrent mutation.
If several worker threads mutate \api{source} at the same time, the source list's
single-owner convention has already been violated. A better design is often:
workers produce immutable results, an EDT task applies those results to the
presentation list, and Swing adapters observe the EDT-confined list. The adapter
exists so thread handoff is explicit, not so every caller can ignore ownership.

\api{ObservableListModel} translates Corusco list changes into Swing
\api{ListDataEvent} deliveries. Inserted ranges become interval-added events.
Removed ranges become interval-removed events. Replacements become
contents-changed events. Moves become contents-changed over the affected range.
Clears become interval removal. The source list remains the data owner; the
Swing model owns only event translation and subscription lifecycle.

\begin{lstlisting}[caption={A JList can display a presenter-owned row source.}]
ObservableListModel<CustomerChoice> model =
        ObservableListModel.of(choices);

JList<CustomerChoice> list = new JList<>(model);
\end{lstlisting}

The list component still owns selection UI, cell rendering, focus, keyboard
behavior, and scrolling. The observable list owns option data. This separation
matters when the same choices also feed a combo box, command enablement, or test.
The model adapter should be closed with the screen because it owns a source-list
subscription.

\api{ObservableComboBoxModel} extends the list model idea for combo boxes. It
routes mutable combo-box operations back to the source list, so user edits and
presenter edits share the same observable state. It owns Swing selection state
separately because Swing permits a selected item that is not present in the
model. If source changes remove the selected value and no equal value remains,
the adapter clears selection and fires the standard selection-change event.

\begin{lstlisting}[caption={Combo-box model selection is not list membership.}]
ObservableComboBoxModel<CustomerChoice> model =
        ObservableComboBoxModel.of(choices);

combo.setModel(model);
model.setSelectedItem(previousChoice);
\end{lstlisting}

This distinction is subtle and practical. A combo box may temporarily select a
typed value that is not in the list while an editor is active, or while options
reload. A source-list removal may invalidate selection. The adapter handles the
Swing contract, while the source list remains the option owner. Presenter code
that cares about selected customer identity should still use a selected-id value
or binding rather than treating combo-box selection as the durable fact.

\api{ObservableTableModel} connects an observable row source with typed table
columns. The row list owns rows. The table descriptor owns column count, column
ids, headers, classes, value extraction, and edit policy. The table model owns
translation into \api{TableModelEvent}. The \api{JTable} owns the viewport,
selection UI, renderers, editors, column model, and row sorter. Keeping these
owners separate is the difference between a table architecture and a large
method named \api{getValueAt}.

\begin{lstlisting}[caption={ObservableTableModel is the adapter, not the row owner.}]
ObservableArrayList<CustomerRow> rows = ObservableArrayList.empty();
TableDescriptor<CustomerRow> descriptor = customerTableDescriptor();
ObservableTableModel<CustomerRow> model =
        ObservableTableModel.of(rows, descriptor);

table.setModel(model);
\end{lstlisting}

Editable columns can replace immutable rows. This is important for modern Java
code that uses records and wither-like methods. A cell editor commits a value.
The column updater creates a replacement row. The observable list emits a
replacement. The table model fires an update. Undo, validation, and diagnostics
can observe row replacement as a data event instead of guessing which bean
property changed.

Table sorting creates another coordinate boundary. If \api{JTable} has a
\api{RowSorter}, view row zero is not model row zero. A table selection binding
must use Swing's conversion APIs when writing selected row state. An observable
list does not remove this Swing rule. It gives the model side precise rows and
events so the binding has something stable to preserve.

\begin{lstlisting}[caption={View rows must become model rows before reading the list.}]
int viewRow = table.getSelectedRow();
CustomerRow row = null;
if (viewRow >= 0) {
    int modelRow = table.convertRowIndexToModel(viewRow);
    row = rows.get(modelRow);
}
selectedCustomer.setValue(row, ChangeOrigin.USER);
\end{lstlisting}

Generated table bindings can hide this boilerplate, but they cannot change the
rule. Selection is meaningful only when the coordinate system is known. A row id
is often better than a row object when refresh may replace row instances. A row
object may be acceptable when row instances are durable and immutable. A model
row index is usually the weakest choice for long-lived state.

Closing adapters is not optional. \api{ObservableListModel},
\api{ObservableComboBoxModel}, \api{ObservableTableModel}, filtered views,
sorted views, mapped views, EDT wrappers, and GlazedLists adapters all own
subscriptions. The Swing component may become undisplayable, but the source list
does not know that unless the adapter is closed or a lifecycle scope closes it.
The user sees a closed window; the heap may still see a subscribed model.

\begin{lstlisting}[caption={A screen closes the adapters it installs.}]
final class CustomerTableView implements AutoCloseable {
    private final ObservableTableModel<CustomerRow> tableModel;

    CustomerTableView(ObservableList<CustomerRow> rows, JTable table) {
        tableModel = ObservableTableModel.of(rows, customerTableDescriptor());
        table.setModel(tableModel);
    }

    @Override
    public void close() {
        tableModel.close();
    }
}
\end{lstlisting}

Later binding scopes can own this cleanup. The chapter shows it manually because
the responsibility should be understood before it is automated. A model adapter
is not merely a value object; it subscribes to a source and talks to Swing.

Renderer performance still matters. Observable lists can make change delivery
precise, but a renderer that performs repository calls, allocates expensive
formatters per cell, or forgets to reset state can still ruin the table. The
list chapter and table chapter meet here: precise row events keep the table from
doing unnecessary work, while disciplined renderers keep visible-cell work
cheap. Neither replaces the other.

Editor behavior matters too. A table cell editor may be active when the row list
changes. A save command should stop or cancel active editing before reading
model state. A row replacement may invalidate the edited row. A source-list
clear should normally cancel editing. These are Swing table workflow rules, but
observable row events make them visible. A broad reset hides what happened; a
precise remove or replace lets the adapter and controller choose a policy.

\section{Large-Data Workflows and Migration Reviews}

Large data does not always mean many rows in memory. It means the screen cannot
afford to be vague. A search screen may show fifty rows from a server result of
millions. A monitoring table may update hundreds of rows per second. A work
queue may contain a few thousand rows but require stable selection and cheap
refresh. A combo box may contain enough options that filtering matters. The
observable-list design should match the workflow, not merely the maximum row
count.

For a repository-backed search, prefer a load state value plus a presentation
list. The load state says idle, loading, loaded, failed, or cancelled. The list
contains the current visible rows. A new search may clear the list, replace it
with a page, or preserve old rows while a busy overlay appears. These are user
experience decisions. The observable list should report whatever row change the
decision implies. The load state should not be smuggled into a fake row.

\begin{lstlisting}[caption={Search state and visible rows are different facts.}]
SimpleValue<SearchState> searchState =
        SimpleValue.of(new SearchState.Idle());
ObservableArrayList<CustomerRow> visibleRows =
        ObservableArrayList.empty();

void publishResults(List<CustomerRow> rows) {
    visibleRows.batch(list -> {
        list.clear();
        rows.forEach(list::add);
    });
    searchState.setValue(new SearchState.Loaded(rows.size()));
}
\end{lstlisting}

For refresh, clearing and repopulating may be the wrong event story. If the user
has selected a row, expanded groups, adjusted columns, and started editing, a
full clear destroys too much context. A refresh can compare old and new rows by
id, replace changed rows, insert new rows, remove missing rows, and preserve
selection by id. This diff can be simple or sophisticated. The important point
is that row identity makes a more truthful update possible.

\begin{lstlisting}[caption={Refresh by id instead of blind reset.}]
void refreshRows(List<CustomerRow> freshRows) {
    Map<CustomerId, CustomerRow> freshById = indexById(freshRows);
    rows.batch(list -> {
        removeRowsMissingFrom(freshById, list);
        replaceRowsThatChanged(freshById, list);
        insertNewRowsInDisplayOrder(freshRows, list);
    });
}
\end{lstlisting}

This listing is a sketch because real diffing depends on order policy and data
volume. It is still a better starting point than \api{clear} followed by
\api{addAll} for every refresh. A row-level diff keeps selection, scrolling, and
rendering calmer. It also produces better diagnostics: the support log can say
"three rows replaced, one inserted, one removed" instead of "table reset."

For streaming updates, batching is a pressure valve. A feed may deliver many
small changes quickly. Applying every update as its own Swing event may keep the
model truthful and still overwhelm the EDT. Coalesce related updates into
batches at a known boundary: timer tick, server message, transaction, or
generation. The batch should still describe the ordered row changes. Do not
solve event pressure by hiding everything behind a reset unless the reset is
actually the truthful operation.

\begin{lstlisting}[caption={Coalesce feed updates into one ordered batch.}]
void applyFeedMessage(List<RowUpdate> updates) {
    rows.batch(list -> {
        for (RowUpdate update : updates) {
            applyOne(update, list);
        }
    });
}
\end{lstlisting}

For paged data, the presentation list may represent only one page. Page changes
can be broad resets because the visible slice really changed. Selection may be
stored outside the page as an id. A detail pane may remain open for a selected
id even when that row is not on the current page. The observable list then owns
page contents, while a scalar value owns selected identity and another state
value owns page/load status. Mixing those facts into the table model makes page
navigation hard to test.

For virtualized data, Swing's ordinary \api{JTable} model can expose a large row
count and load cells on demand, but that is a different contract from an
in-memory \api{ObservableArrayList}. Corusco observable lists are excellent for
presentation collections and eventful row caches. They are not a demand-paged
database cursor by themselves. If the application needs true virtualization,
design a virtual table model or a paged source, then use observable lists for
the slices or caches that are actually present.

GlazedLists interop is valuable in this part of the design. GlazedLists already
has mature event-list machinery for filtering, sorting, and Swing integration.
\api{GlazedObservableList} adapts an \api{EventList} into Corusco's
\api{ObservableList} contract. The wrapped event list remains the storage owner.
The adapter uses the event list's locks for reads and direct adapter mutations.
Closing the adapter removes its listener but does not dispose the wrapped list.
That ownership sentence should appear in code reviews whenever an adapter is
introduced.

\begin{lstlisting}[caption={Adapt GlazedLists deliberately at the boundary.}]
EventList<CustomerRow> eventRows = new BasicEventList<>();
GlazedObservableList<CustomerRow> rows =
        GlazedObservableList.of(eventRows);

try {
    installCoruscoTable(rows);
} finally {
    rows.close();      // removes adapter listener
}
\end{lstlisting}

Use GlazedLists when it should own the event-list pipeline. Use Corusco core
lists when the presenter owns simple row storage or projections. Do not stack
several sorting and filtering systems casually. A table row sorter over a
Corusco sorted view over a GlazedLists sorted list is probably a design smell
unless each layer has a clearly different purpose.

Migration from hand-written \api{AbstractTableModel} should begin with event
truth, not code generation. Find every \api{fireTableDataChanged}. Ask whether
the model knows a more precise change. Replace the broad reset with a row insert,
delete, update, or ordered batch where possible. Then extract row storage into an
observable list. Then introduce typed column descriptors. Then let a Corusco
table model translate events. This order preserves behavior while reducing
custom code.

\begin{lstlisting}[caption={A staged migration from a custom table model.}]
// Step 1: make the existing model precise.
fireTableRowsUpdated(index, index);

// Step 2: move row storage into an observable source.
rows.set(index, updatedRow);

// Step 3: let the adapter translate the row event.
table.setModel(ObservableTableModel.of(rows, descriptor));
\end{lstlisting}

Migration from a \api{DefaultListModel} or \api{DefaultComboBoxModel} follows the
same idea. First identify who owns the options. If the component model is the
only owner and the list is small, the default model may be enough. If presenter
state, tests, commands, or other components need the same options, move the data
into an observable list and adapt it. The component should present options; it
should not be the only repository of application state.

Tests should cover both model events and adapter events. Core tests mutate an
\api{ObservableArrayList}, filtered view, sorted view, mapped view, or loadable
list and assert \api{ListChangeSet} contents. Swing tests construct
\api{ObservableListModel}, \api{ObservableComboBoxModel}, or
\api{ObservableTableModel} on the EDT and assert Swing event ranges. The first
kind of test proves collection semantics. The second proves Swing translation.
Do not rely on a visual table test to prove every list event rule.

\begin{lstlisting}[caption={Core test for precise list changes.}]
ObservableArrayList<String> rows = ObservableArrayList.of(List.of("b", "c"));
List<ListChangeSet<String>> events = new ArrayList<>();
rows.subscribe(events::add);

rows.add(0, "a");
rows.set(1, "B");
rows.move(2, 0);
rows.remove(1);

assertThat(events).hasSize(4);
assertThat(events.getFirst().origin()).isEqualTo(StandardChangeOrigin.MODEL);
\end{lstlisting}

Adapter tests should assert EDT delivery when the adapter claims it. An
\api{EdtObservableList} test can mutate the source from a worker, wait for a
callback, and assert that the callback ran on the EDT. An \api{ObservableListModel}
test can record \api{ListDataEvent} ranges. An \api{ObservableTableModel} test
can record \api{TableModelEvent} types and rows. These tests are small and catch
the boundary failures that otherwise appear as intermittent Swing bugs.

Performance tests should use realistic row counts and realistic row shapes. A
five-row table proves wiring. It does not prove refresh behavior, renderer cost,
memory footprint, or selection preservation. A useful large-data smoke test
loads thousands of deterministic rows, applies a filter, replaces a row, moves a
row, preserves selection by id, and records the event count. It need not render
a window. The table screenshot can remain modest; the model test should carry
the scale.

Diagnostics should summarize list operations. A debug status line or support
dump can report source row count, visible row count, last change set size, last
event type, selected id, filter text, sort policy, and adapter thread policy.
When a user reports that a table jumped, these facts matter. Without them, the
developer is left guessing whether the cause was reset, sorter, selection, data
load, or repaint.

\begin{lstlisting}[caption={A small list diagnostic snapshot.}]
record RowDiagnostics(
        int sourceRows,
        int visibleRows,
        int lastChangeCount,
        String sortPolicy,
        String filterText,
        Object selectedId,
        boolean deliveredOnEdt) {
}
\end{lstlisting}

Do not let diagnostics become another observer leak. If a diagnostics panel
subscribes to row sources, close it. If a support bundle samples snapshots, take
immutable snapshots and release them. Large-data diagnostics can accidentally
become large-data retention bugs when they keep old row lists for comparison
without a retention policy.

Observable lists also affect accessibility. A list or table that receives
precise events can preserve focus and accessible context more gracefully than a
table that resets on every update. If the selected row remains the same domain
object, accessible descriptions can remain stable. If a refresh removes the row,
the screen can announce a meaningful status. This is another reason row identity
is not merely a persistence concern.

The final migration review is simple. Ask whether the table tells the truth
about row changes. Ask whether rows have stable identity where continuity
matters. Ask whether filtering and sorting have one owner. Ask whether Swing sees
events on the EDT. Ask whether adapters are closed. Ask whether tests cover the
core list and the Swing adapter separately. If the answers are yes, the row
collection is ready to support richer Corusco table, form, and command chapters.

\section{Worked Row-Source Scenarios}

The first scenario is the search result table. The user enters criteria, presses
Search, and receives a bounded result set. The screen usually needs a load state,
a visible row list, selected identity, status text, and commands for refresh,
open, export, and clear. The row list should not carry all of those facts. It
should carry the visible rows. The load state value says whether the query is
idle, running, successful, empty, failed, or cancelled. The selected-id value
says which row is current. Derived values connect those facts to commands.

\begin{lstlisting}[caption={Search screen state around a visible row list.}]
ObservableArrayList<CustomerRow> visibleRows =
        ObservableArrayList.empty();
SimpleValue<CustomerId> selectedId = SimpleValue.empty();
SimpleValue<SearchState> searchState =
        SimpleValue.of(new SearchState.Idle());
CollectionDerivedValue<CustomerRow, Boolean> rowsEmpty =
        CollectionDerivedValue.isEmpty(visibleRows);

DerivedValue<Boolean> canExport = DerivedValue.of(
        () -> !rowsEmpty.value()
                && searchState.value() instanceof SearchState.Loaded,
        rowsEmpty, searchState);
\end{lstlisting}

For predicate-style questions, use the stream-based derived factories:

\begin{lstlisting}[caption={Stream-based collection derivation.}]
CollectionDerivedValue<CustomerRow, Boolean> hasInvalidRows =
        CollectionDerivedValue.anyMatch(visibleRows,
                CustomerRow::hasValidationProblems);
\end{lstlisting}

\begin{lstlisting}[caption={Expose row count as scalar state when commands need it.}]
SimpleValue<Integer> rowCount = SimpleValue.of(visibleRows.size());
Subscription rowCountSubscription = visibleRows.subscribe(changes ->
        rowCount.setValue(visibleRows.size(), ChangeOrigin.MODEL));

DerivedValue<Boolean> canExport = DerivedValue.of(
        () -> rowCount.value() > 0
                && searchState.value() instanceof SearchState.Loaded,
        rowCount, searchState);
\end{lstlisting}

This pattern is small and useful. A table model observes row changes for Swing.
A command observes row count as scalar state. A status label can say "34 rows."
Tests can mutate the list and assert command enablement. The row-count
subscription must still be owned by a scope or presenter. There is no free
lifecycle merely because the state is derived from a list.

The second scenario is refresh with selection preservation. Suppose a user
selects customer 42. A refresh arrives with updated rows. The selected row may
still exist, may have new values, may be filtered out, or may have disappeared.
The table should preserve selection when the id survives. It should clear
selection when the id disappears. It should avoid jumping when unrelated rows
change. This requires row identity before it requires a table API.

\begin{lstlisting}[caption={Preserve selected id across a refresh.}]
void afterRefresh() {
    CustomerId selected = selectedId.value();
    if (selected == null) {
        return;
    }
    boolean stillPresent = rows.stream()
            .anyMatch(row -> row.id().equals(selected));
    if (!stillPresent) {
        selectedId.setValue(null, ChangeOrigin.MODEL);
    }
}
\end{lstlisting}

The table binding can then select the visible row matching \api{selectedId} if
one exists. If a filter hides the row, the selected id may remain a valid
semantic selection even though no visible table row is selected, or the screen
may choose to clear it. That is a workflow decision. The important point is that
the decision is no longer a side effect of row indexes.

The third scenario is incremental validation in a row list. A table of imported
records may show row-level problems. The row object can include a problem
summary, or a separate problem map can be keyed by row id. Replacing one row
should update one row. Recomputing validation for every row after one edit may
be necessary if rules are global, but it should be a named validation pass, not
a renderer trick.

\begin{lstlisting}[caption={Row problems can be data, not renderer state.}]
record ImportedCustomerRow(
        CustomerId id,
        String displayName,
        ProblemSet problems) {
}

void replaceProblems(CustomerId id, ProblemSet problems) {
    int index = findRowIndex(id);
    ImportedCustomerRow old = rows.get(index);
    rows.set(index, new ImportedCustomerRow(
            old.id(), old.displayName(), problems));
}
\end{lstlisting}

This design lets the renderer present problem severity cheaply. It also lets a
problem summary observe the row list or a separate problem model. The renderer
does not validate. The table model does not own problems. The row source
publishes replacement when row presentation facts change.

The fourth scenario is a combo box with mutable options. A user may type a new
tag, insert it into the option list, and select it. Swing's
\api{MutableComboBoxModel} supports option mutations. Corusco's combo-box
adapter routes those mutations back to the source list. This is useful only when
the option source is allowed to change from the combo. If options come from a
controlled vocabulary, the combo should not expose arbitrary mutation. The source
ownership decision comes before the adapter choice.

\begin{lstlisting}[caption={Combo mutations are source-list mutations.}]
ObservableArrayList<TagChoice> tags = ObservableArrayList.empty();
ObservableComboBoxModel<TagChoice> model =
        ObservableComboBoxModel.of(tags);

model.addElement(new TagChoice("urgent"));
assert tags.stream().anyMatch(new TagChoice("urgent")::equals);
\end{lstlisting}

Selection still deserves its own model when other parts of the screen care. The
combo-box model selection is a Swing fact. A presenter value such as
\api{selectedTag} is a screen fact. A binding can connect them. This distinction
keeps commands and tests from reaching into a component model.

The fifth scenario is a master table with a detail list. Selecting a project
loads tasks. The project list and task list are separate row sources. The
selected project id is a scalar value. The task list may be loadable or task
backed. A project refresh should preserve selected project id if possible. A
task refresh should not reset project selection. These relationships are easy to
blur when everything is written inside one table listener.

\begin{lstlisting}[caption={Master id drives a separate detail row source.}]
SimpleValue<ProjectId> selectedProjectId = SimpleValue.empty();
ObservableArrayList<TaskRow> taskRows = ObservableArrayList.empty();

selectedProjectId.subscribe(event -> {
    taskRows.clear();
    if (event.newValue() != null) {
        loadTasksFor(event.newValue());
    }
});
\end{lstlisting}

The listing is deliberately plain. In production, loading may be a task with
cancellation and stale-result suppression. The chapter's point is ownership:
project selection is a value; tasks are a list; loading is a task; the detail
table is an adapter. A single table listener should not own all four.

The sixth scenario is a live monitoring table. Rows update frequently. The user
may sort by severity or timestamp. The screen must avoid both per-cell chaos and
broad resets. A common strategy is to maintain a map by id for incoming updates,
then apply changes to the observable list in timed batches on the EDT. If sort
order is view-local, let the table sorter handle it. If sort order is part of
the presentation model, publish a sorted view or use an event-list engine.

\begin{lstlisting}[caption={Apply live updates on a scheduled presentation boundary.}]
void flushPendingUpdates() {
    List<RowUpdate> updates = pending.drain();
    if (updates.isEmpty()) {
        return;
    }
    rows.batch(list -> {
        for (RowUpdate update : updates) {
            applyUpdateById(update, list);
        }
    });
}
\end{lstlisting}

The scheduled boundary can be a Swing timer, a task callback, or a domain
message boundary. What matters is that the UI receives coherent updates at a
rate it can render. The observable list is not a rate limiter by itself. It is
the mechanism that describes the mutations once the application chooses when to
publish them.

The seventh scenario is failure. A row load can fail. A filter can be invalid.
A repository can time out. A GlazedLists adapter can be closed. A background
update can arrive after the view is gone. The list should not encode failure by
inserting a fake row named "Error". Use a load or problem state. Keep rows as
rows. Then the table can show an overlay, status line, or problem summary
without corrupting the row model.

\begin{lstlisting}[caption={Failure is state beside the rows, not a fake row.}]
ObservableArrayList<CustomerRow> rows = ObservableArrayList.empty();
SimpleValue<ProblemSet> loadProblems =
        SimpleValue.of(ProblemSet.empty());

void loadFailed(Throwable failure) {
    loadProblems.setValue(problemFor(failure), ChangeOrigin.SYSTEM);
}
\end{lstlisting}

This matters for export and automation. A fake error row might be exported as a
customer. It might be selected, edited, counted, or sent to a detail loader.
Modeled failure state can be presented visually without pretending to be data.

The eighth scenario is undo. If table edits are undoable, row replacement events
are helpful. An undoable edit can store old row and new row, or a domain command
that knows how to reapply the change. The observable list reports the visible
replacement. Undo should not be implemented by asking the table what cell was
last painted. The row source is the right boundary.

\begin{lstlisting}[caption={Undoable row replacement records model facts.}]
record ReplaceRowEdit(int index, CustomerRow oldRow, CustomerRow newRow) {
    void redo(ObservableList<CustomerRow> rows) {
        rows.set(index, newRow);
    }

    void undo(ObservableList<CustomerRow> rows) {
        rows.set(index, oldRow);
    }
}
\end{lstlisting}

For sorted or filtered views, storing only the visible index may be too weak.
Store row identity and enough context to find the row in the source. Again,
coordinate systems matter. Undo belongs to the model operation, not to whatever
view position happened to display the row at edit time.

The ninth scenario is generated descriptors. A generated table descriptor can
name columns, value extractors, editors, resource keys, and perhaps row identity.
It can build an observable table model from an observable row source. It should
not own the row source itself unless the generated companion is explicitly a
screen model. Generation is strongest when handwritten code already separates
row ownership, column meaning, and Swing presentation.

The tenth scenario is screenshot production. A book screenshot should use a
deterministic row source with seeded data and stable sorting. It should avoid
loading from a real repository during capture. It should set selection, filter,
busy state, and problem state explicitly. Observable lists make this easy:
construct the row source, apply known mutations, install the adapter, capture
the component. If a screenshot requires five seconds of timing luck, it is not a
good companion example.

These scenarios all reduce to one review question: what is the collection fact,
and who owns it? Once that answer is clear, the rest of the design becomes
ordinary. A filter is a view over a source. A sort is either a view preference or
a presenter-owned order. A table model is an adapter. A combo selection is a
Swing fact unless bound to a presenter value. A large repository result is not
automatically an in-memory list. A row event should say what changed. This is the
discipline that makes Swing large-data screens feel calm.

For chapter acceptance, a row-source example should prove more than "the table
has rows." It should show an insertion, a replacement, a removal, and a refresh.
It should preserve selection by identity where the workflow promises continuity.
It should run at least one model-level test without Swing. If it has a Swing
adapter, it should run one EDT test that records Swing event ranges. If it has a
screenshot, the data should be deterministic and the table should not depend on
external services during capture. These checks are small enough for examples and
strong enough to keep the prose honest.

For code review, start with the row source and walk outward. The row source owns
membership and order. Derived views own filtered, sorted, or mapped projections.
Adapters own Swing event translation. Components own interaction and painting.
Commands own operations. Values own scalar facts such as selection, row count,
filter text, load state, and disabled reason. When a class owns too many of
these at once, it becomes hard to test and harder to close.

For performance review, ask what happens when the row count grows by two orders
of magnitude. Does the refresh still diff by identity or does it reset? Does the
filter run in a renderer? Does every update allocate expensive presentation
objects? Does selection survive? Does the adapter receive events on the EDT? Are
events coalesced at a sensible boundary? Does the screenshot harness use a
realistic number of rows or only a toy table? These questions are not premature
optimization. They are how one prevents the classic Swing table that works in a
demo and collapses under production data.

For lifecycle review, ask which objects subscribe and who closes them. A
filtered view subscribes to its source. A sorted view subscribes to its source. A
mapped view subscribes to its source. A Swing list, combo, or table model
subscribes to its source. A diagnostics panel may subscribe too. The source list
will not know that a window closed unless those subscriptions are closed. Large
data leaks are especially painful because the retained objects are large enough
to be noticed but often hidden behind innocent adapters.

For user-experience review, ask what the user sees during change. A precise row
replacement should not collapse selection. A refresh should not jump to the top
unless the result truly changed that way. A filter change may reset the visible
projection, but the screen should make that policy understandable. A failure
should be shown as state, not as a fake data row. A busy load should not freeze
the table. Observable lists do not solve these UX questions automatically; they
give the program enough truthful events to answer them cleanly.

For architecture review, ask whether the list chapter's lessons are being
carried into later Corusco features. A Corusco table descriptor without stable
row identity cannot preserve selection well. A binding scope that installs a
table model but never closes it leaks. A background task that publishes a list
from a worker thread without an EDT boundary violates Swing. A generated table
companion that knows columns but not row ownership is only half the story. The
observable-list contract is the collection foundation beneath those later
features.

The next chapters use lists as ordinary building blocks. A form may edit the
selected row from a list. A command may be enabled only when a list has selected
rows. A background task may replace a row source after loading. A generated table
may create descriptors over rows whose membership is owned elsewhere. A real
screen may combine all of these. If this chapter has done its job, the reader
will not see a table as a special widget full of mysterious callbacks. The
reader will see a row source, a set of projections, a table adapter, and a few
scalar values that describe selection and workflow.

That mental model is intentionally humble. It does not promise that every data
problem is small. It does promise that when a problem becomes large, the program
has named places to put the complexity: repository, paging source, event-list
adapter, observable presentation list, descriptor, binding scope, task service,
and diagnostics. Swing remains Swing, but the row story is no longer hidden
inside a table subclass.

In practice this is what makes a large Swing table pleasant to maintain: the
boring parts stay boring. Adding a row is a list insertion. Replacing a row is a
list replacement. A filter is a projection. A table model is an adapter. A
selection is a value. A refresh is a published set of mutations, not a panic
button. Once those sentences are true, the screen has room for harder domain
problems without turning every table bug into archaeology.
The collection model becomes a map, not a maze. That map is what later chapters
reuse.

\section{Exercises}

\begin{exercisebox}
\begin{enumerate}
\item Replace one \api{fireTableDataChanged} call with a precise inserted,
removed, or updated range.
\item Write a test for an \api{ObservableArrayList} batch and assert the order
of delivered changes.
\item Create a filtered row view and verify that changing a source row can add
or remove it from the filtered view.
\item Adapt a source list with \api{EdtObservableList}. Confirm that Swing
subscribers run on the EDT.
\item Compare a Swing \api{RowSorter} with a presenter-owned
\api{SortedList}. Decide which one owns a user-selected sort in your screen.
\item Add stable row identity to a refresh path and preserve selection across
row replacement.
\end{enumerate}
\end{exercisebox}

\textbf{Checklist.}

\begin{itemize}
\item Fire precise list changes; avoid broad resets as a habit.
\item Use batches for related bulk mutations.
\item Keep filtering, sorting, and mapping out of renderers.
\item Marshal list changes before they reach Swing models.
\item Close list subscriptions and table models with the screen.
\item Test large-data behavior with thousands of rows, not five.
\item Preserve row identity when selections must survive refresh.
\item Use GlazedLists interop deliberately when its event-list machinery is the
right owner.
\end{itemize}
